\documentclass[11pt]{article}
\usepackage{amsmath, amsthm, amssymb, amsfonts}
\usepackage{fullpage}
\usepackage{hyperref}

\title{The Proof of the Mass Gap in Yang-Mills Theory}
\author{[Your Name]}
\date{\today}

\begin{document}

\maketitle

\section{Introduction}
This document presents the proof of the mass gap in Yang-Mills theory, an essential problem in mathematical physics. The proof is structured into three key steps: establishing the existence of a unique vacuum state, demonstrating that all excited states have a positive energy above the vacuum, and showing that the Wilson loop behavior implies confinement, thereby proving the mass gap.

\section{The Proof of the Mass Gap}

\subsection{Step 1: Existence of a Unique Vacuum State}

The first step in proving the mass gap is to establish the existence of a unique vacuum state \( |0\rangle \) with the lowest possible energy. The Yang-Mills Hamiltonian \( H_{\text{YM}} \) is derived from the Yang-Mills Lagrangian and governs the dynamics of the gauge fields. The vacuum state \( |0\rangle \) is defined as the state that minimizes the expectation value of the Hamiltonian:
\[
E_0 = \langle 0 | H_{\text{YM}} | 0 \rangle
\]
This state corresponds to a pure gauge configuration where the field strength \( F_A = 0 \). The uniqueness of the vacuum state follows from the compactness of the gauge group \( G \) and the positivity of the Hamiltonian. Since the vacuum is the lowest energy state, any perturbation would increase the energy, thus reinforcing the uniqueness.

\subsection{Step 2: Positive Energy for Excited States}

The second step involves showing that any state \( |E_1\rangle \) orthogonal to the vacuum state \( |0\rangle \) must have a positive energy above the vacuum. This step is critical because it establishes the mass gap.

The Hamiltonian \( H_{\text{YM}} \) has a spectrum of eigenvalues corresponding to the energy levels of the theory. The vacuum energy \( E_0 \) is the lowest eigenvalue. For any excited state \( |E_1\rangle \), the energy difference relative to the vacuum is given by:
\[
E_1 - E_0 \geq m > 0
\]
This inequality holds because the spectrum of \( H_{\text{YM}} \) is discrete and bounded below, a property that arises from the compactness of the gauge group \( G \). The existence of a positive gap \( m \) ensures that the theory does not admit massless excitations other than the gauge bosons associated with \( G \).

\subsection{Step 3: Wilson Loop and Confinement Implies Mass Gap}

The final step in the proof relates the mass gap to the behavior of the Wilson loop. The Wilson loop \( W(C) \) is a gauge-invariant observable defined for a closed loop \( C \) in spacetime:
\[
W(C) = \text{tr}\left( P \exp \left( \oint_C A \right) \right)
\]
The area law for the Wilson loop, expressed as:
\[
\langle W(C) \rangle \sim \exp(-\sigma \cdot \text{Area}(C)),
\]
indicates that the potential between quarks increases linearly with distance, leading to confinement. The parameter \( \sigma \) is the string tension, and its non-zero value implies that quarks are confined within hadrons.

This confinement behavior suggests that there are no massless particles (except the gauge bosons themselves) and that the correlation functions decay exponentially with distance. The exponential decay is characterized by a correlation length \( \xi \), which is related to the mass gap \( m \) by:
\[
m \sim \frac{1}{\xi}
\]
Thus, the area law behavior of the Wilson loop implies a finite correlation length and consequently a positive mass gap in the theory.

\section{Conclusion}

The proof of the mass gap in Yang-Mills theory has been established through three critical steps: the existence of a unique vacuum state, the positive energy of all excited states, and the Wilson loop's area law behavior indicating confinement. These steps together provide a rigorous foundation for the mass gap, ensuring that Yang-Mills theory with a compact, simple Lie group \( G \) possesses a positive mass gap.

This result has profound implications for our understanding of non-perturbative quantum field theory, particularly in quantum chromodynamics (QCD), where the mass gap underpins the confinement of quarks and the stability of hadronic matter.

\end{document}